\documentclass[10pt,a4paper]{amsart}
\usepackage[margin=19mm]{geometry}
\usepackage{amsmath,amssymb,mathtools}
\usepackage[scr=rsfs,cal=euler]{mathalfa}
\usepackage{xfrac}
\usepackage[unicode,hidelinks,bookmarks=false]{hyperref}
\newcommand{\ie}{i.e.\ }
\newcommand{\cf}{cf.\ }
\newcommand{\resp}{respectively\ }
\newcommand{\Subsection}{\subsection}
\providecommand{\nobreakdash}{\nobreak\hbox{-}}
\setlength{\emergencystretch}{2em}
\multlinegap0pt
\usepackage{bm}
\usepackage{bbm}
\usepackage{dsfont}
\usepackage{xspace}
\usepackage{cancel}
\usepackage{mathtools}
\usepackage{amsthm}
\makeatletter
\@ifundefined{Theorem}{\newtheorem{Theorem}{Theorem}}{}
\@ifundefined{Corollary}{\newtheorem{Corollary}[Theorem]{Corollary}}{}
\@ifundefined{Definition}{\newtheorem{Definition}[Theorem]{Definition}}{}
\@ifundefined{Lemma}{\newtheorem{Lemma}[Theorem]{Lemma}}{}
\@ifundefined{Proposition}{\newtheorem{Proposition}[Theorem]{Proposition}}{}
\@ifundefined{Remark}{\newtheorem{Remark}[Theorem]{Remark}}{}
\makeatother
\title[On the center of distances of finite ultrametric spaces]{On the center of distances of finite ultrametric spaces}
\author{Oleksiy Dovgoshey, Olga Rovenska}
\date{Source: arXiv:2603.25850v1 (CC BY 4.0)}
\begin{document}
\maketitle

\noindent\emph{Source and licence.} On the center of distances of finite ultrametric spaces. Version arXiv:2603.25850v1 (CC BY 4.0), distributed under the Creative Commons Attribution 4.0 International licence, as recorded in the arXiv licence metadata for this exact version and shown on the abstract page https://arxiv.org/abs/2603.25850v1. Full rights notice: licenses/arxiv-2603.25850-CC-BY-4.0.txt.\par\smallskip

\noindent\textit{Editorial scope.} Counted unit: the Proposition whose source label is \texttt{woorf} in the article (source0.tex lines 982--993, source label \texttt{woorf}), delivered together with the article's own complete original proof of it (source0.tex lines 995--1267, one \texttt{proof} environment of 273 lines). No mathematics is written, repaired or re-derived: statement and proof are the article's own text line for line. Required context retained uncounted: jgds (lines 307--321); zbh1 (lines 340--347); saqh (lines 361--372); mnbffss (lines 375--376); edftt (lines 651--665); ugrdd (lines 715--721); sefugi (lines 728--740); sabra (lines 924--935). Every definition, display and label of those lines is retained unchanged. Omitted from this package and not redistributed: 1--306, 322--339, 348--360, 373--374, 377--650, 666--714, 722--727, 741--923, 936--981, 994--994, 1268--1861. Nothing inside a retained excerpt is omitted or altered: every retained body is delivered exactly as the article's lines are written, so no omission receipt of any kind accompanies this package. Citations: the retained excerpts carry no citation command at all, so no bibliography entry of the article is transcribed and no reference is redistributed. Reference adaptation: none was needed and none was made; every reference inside the delivered unit resolves inside it, so no unresolved reference or citation is produced. Printed slips kept literal: no printed word, symbol, hypothesis or number of the counted statement or of its proof was altered. Layout and class: the article's own class is \texttt{sn-jnl} with packages including graphicx, multirow, amsmath, amssymb, amsfonts, amsthm, mathrsfs, appendix, xcolor, textcomp, manyfoot, booktabs, algorithm, algorithmicx; the wrapper is the lane's pinned \texttt{amsart} editorial wrapper at 10pt on A4, which restates the theorem-like environments the retained text needs and adds the article's own macro definitions that the retained text needs (none), with extra wrapper packages (bm, bbm, dsfont, xspace, cancel, mathtools). The delivered numbering is the unit's own. Assets: no retained span contains a figure environment, a graphics-inclusion command, a PDF-page inclusion, a table or a TikZ picture; no image file is copied into this package, the package declares no asset dependency and no image file is redistributed with it. Licence: version arXiv:2603.25850v1 of this article is distributed under the Creative Commons Attribution 4.0 International licence, as recorded in the arXiv licence metadata for this exact version and shown on the abstract page https://arxiv.org/abs/2603.25850v1. A full rights notice is delivered at licenses/arxiv-2603.25850-CC-BY-4.0.txt. Characters: every retained excerpt is pure ASCII, so the delivered engine, XeLaTeX with its default Latin Modern text font, reports no missing character.
\begin{Proposition}
\label{jgds}
Let $G$ be a connected finite graph with $|V(G)|\geq 2$. Then the following conditions are equivalent.
\begin{enumerate}
    \item[(i)] $G$ is a tree.

        \item[(ii)] The equality 
\begin{equation*}
    |V(G)|=1+|E(G)|
\end{equation*}
holds.

            \item[(iii)] Any two distinct vertices of $G$ are joined by a unique path in $G$.
\end{enumerate}
\end{Proposition}

\begin{equation}
    \label{zbh1}
\delta_T^{+}(v):=
\begin{cases}
\delta_T(v), & \text{if } v=r,\\
\delta_T(v)-1, & \text{if } v\ne r,
\end{cases}
\end{equation}

\begin{Definition}\label{saqh}
Let $T=T(r)$ be a rooted tree and let $u,v$ be distinct vertices of $T$. 
The vertex $u$ is called the {\it direct successor} of $v$ if
\(
u \in V(T)\setminus\{r\},\) \( \{u,v\}\in E(T),
\)
and
\(
v \in V(P_{u,r}),
\)
where $P_{u,r}$ is the path joining the vertex $u$ with the root $r$ in $T$.
\end{Definition}

\begin{Remark}\label{mnbffss} Let $T=T(r)$ be a rooted tree. Condition $(iii)$ of Proposition \ref{jgds} shows that for each vertex $u\in V(T)\setminus\{r\}$ there is the unique vertex $v\in V(T)$ such that $u$ is a direct successor of $v$.
\end{Remark}

\begin{Theorem}\label{edftt}
Let $T = T(r,l)$ be a finite labeled rooted tree with $|V(T)|\geq 3$. Then the following two conditions are equivalent.
\begin{enumerate}
\item[(i)] For every $v \in V(T)$ we have $\delta_T^{+}(v) \ne 1$, the logical equivalence
\[
(\delta_T^{+}(v)=0) \iff (l(v)=0)
\]
is valid and, in addition, the inequality
\(
l(u) < l(v)
\)
holds whenever $u$ is a direct successor of $v$.
\item[(ii)] There exists a finite  ultrametric space $(X,d)$ with $|X|\geq 2$ such that $T(r,l)$ and the representing tree $T_X=T_X(r_X,l_X)$ are isomorphic as labeled rooted trees.
\end{enumerate}
\end{Theorem}

\begin{Proposition}\label{ugrdd}
Let $(X,d)$ be a metric space and let $C(X)$ be the center of distances of $(X,d)$. Then the equality 
\begin{equation*}
C(X) = \bigcap_{p \in X} D_p(X) 
\end{equation*}
holds.
\end{Proposition}

\begin{Corollary}\label{sefugi}
Let $(X,d)$ be a finite  ultrametric space with $|X|\geq 2$ and let 
\(
T_X=T_X(r_X,l_X)
\)
be the representing tree of \((X,d)\).
Then for every $p\in X$  we have the equality
\begin{equation}
    \label{imgt}
D_p(X)=\{l_X(u) : u \in V(P_{\{p\},r_X})\},
\end{equation}
where \(P_{\{p\},r_X}\) is the path in \(T_X\) joining the leaf \(\{p\}\) with the root \(r_X\).
\end{Corollary}

\begin{Lemma}\label{sabra}
Let $(Y,\rho)$ be a finite ultrametric space with $|Y|\ge 2$ and let $T_Y=T_Y(r_Y,l_Y)$ be the representing tree of $(Y,\rho)$. If $p_1,p_2$ are distinct points of $Y$ and $v_0$ is an internal node of $T_Y$ such that $\{p_1\},\{p_2\}$ are direct successors of $v_0$, then the equality
\begin{equation}
\label{kar1}
D_{p_1}(Y)=D_{p_2}(Y)
\end{equation}
holds, where
\[
D_{p_i}(Y):=\{\rho(p_i,y):y\in Y\} 
\]
for $i=1,2.$
\end{Lemma}

\begin{Proposition}\label{woorf}
Let $(X,d)$ be a finite ultrametric space with $|X|\ge 2$. Then there exists a finite ultrametric space $(Y,\rho)$ such that
\begin{equation}
    \label{ned}
|Y|=1+|X|
\end{equation}
and  the equality
\begin{equation}\label{gaf1}
C(X)=C(Y)
\end{equation}
holds, where $C(X)$ and $C(Y)$ are the centers of distances of $(X,d)$ and $(Y,\rho)$, respectively.
\end{Proposition}

\begin{proof}
Let $T_X=T_X(r_X,l_X)$ be the representing tree of $(X,d)$. Since $(X,d)$ is finite and $|X|\ge2$ holds, we can find $x_1 \in X$ and an internal node $v_0$ of $T_X$ such that $\{x_1\}$ is a direct successor of $v_0$ (see Remark~\ref{mnbffss}).

Let us define a graph $T^*$ by equalities
\begin{equation}
    \label{avt1}
V(T^*) := V(T_X)\cup\{v^*\},
\end{equation}
\begin{equation}
    \label{avt2}
E(T^*) := E(T_X)\cup\{v^*,v_0\},
\end{equation}
where $v^*$ is an arbitrary point such that $v^*\notin V(T_X)$.

It is clear that $T^*$ is a connected graph satisfying the equalities
\[
|V(T^*)| = 1 + |V(T_X)|,
\]
\[
|E(T^*)| = 1 + |E(T_X)|.
\]
Hence $T^*$ is a tree by Proposition \ref{jgds}.





Now we determine a root $r^*$ in the tree $T^*$ as
\begin{equation}
    \label{a1}
r^* := r_X,
\end{equation}
where $r_X$ is the root of $T_X$, and define a labeling
\(
l^* : V(T^*) \to \mathbb{R}^+
\)
by the formula
\begin{equation}
    \label{a2}
l^*(u) :=
\begin{cases}
l_X(u), & \text{if } u \in V(T_X),\\
0, & \text{if } u = v^*.
\end{cases}
\end{equation}
Equalities \eqref{avt1}, \eqref{avt2} imply the equalities
\begin{equation}
    \label{be1}
\delta_{T^*}(v^*) = 1,\quad
\delta_{T^*}(v_0) = 1 + \delta_{T_X}(v_0) 
\end{equation}
and the equality
\begin{equation}
    \label{be2}
\delta_{T_X}(v) = \delta_{T^*}(v)
\end{equation}
for each $v \in V(T_X)\setminus\{v_0\}$.
Hence, using \eqref{zbh1} we can rewrite \eqref{be1}, \eqref{be2} as
\begin{equation}
    \label{be3}
\delta^{+}_{T^*}(v^*) = 0,
\quad
\delta^{+}_{T^*}(v_0) = 1 + \delta^{+}_{T_X}(v_0), 
\end{equation}
and
\begin{equation}
    \label{be4}
\delta^{+}_{T_X}(v) = \delta^{+}_{T^*}(v)
\end{equation}
for each $v \in V\setminus\{v_0\}$.


Let us prove that condition $(i)$ of Theorem \ref{edftt} is satisfied for $T(r,l)=T^*(r^*,l^*)$.




Since \(T_X(r_X,l_X)\) is the representing tree of \((X,d)\), condition $(i)$ of Theorem~\ref{edftt} and formulas \eqref{a2}, \eqref{be3},  \eqref{be4} show that the relation
\begin{equation}
    \label{be5}
\delta^{+}_{T^*}(v) \neq 1 
\end{equation}
and the logical equivalence
\begin{equation}
    \label{be6}
(\delta^{+}_{T^*}(v)=0)\;\Longleftrightarrow\; (l^{*}(v)=0) 
\end{equation}
are valid.
Thus it suffices to show that the inequality
\begin{equation}\label{eq:star1}
l^{*}(u)<l^{*}(v)
\end{equation}
holds whenever \(u,v\in V(T^{*})\) and \(u\) is a direct successor of \(v\) in the rooted tree \(T^{*}\).

Let \(u\) be the direct successor of \(v\) in \(T^{*}\). Suppose first that \(u,v\in T_X\).
Since \(T_X\) is a subtree of \(T^{*}\) and \[r^{*}=r_X\] holds by \eqref{a1}, Definition~\ref{saqh} implies that \(u\) is the direct successor of \(v\) in \(T_X(r_X)\).
The tree \(T_X=T_X(r_X,l_X)\) is the representing tree of \((X,d)\).
Hence we have the inequality
\[
l_X(u)<l_X(v)
\]
by Theorem~\ref{edftt}. 
This implies \eqref{eq:star1} by \eqref{a2}.


Let us consider the case when the set \(\{u,v\}\) is not a subset of \(V(T_X)\),
\begin{equation}\label{eq:star2}
\{u,v\}\not\subseteq V(T_X).
\end{equation}

Then using \eqref{avt1}, \eqref{avt2} and Definition~\ref{saqh} we obtain the equality
\begin{equation}\label{eq:star3}
\{u,v\}=\{v^{*},v_0\}.
\end{equation}
Since \(v_0\) is a vertex of \(T_X\), formulas \eqref{eq:star2}, \eqref{eq:star3} show that
\begin{equation}\label{eq:star4}
u=v^{*}\quad\text{and}\quad v=v_0.
\end{equation}
Consequently we have the equalities
\begin{equation}\label{eq:star5}
l^{*}(u)=0
\qquad\text{and}\qquad
l^{*}(v)=l_X(v_0)
\end{equation}
by formula \eqref{a2}.
Moreover, since \(\{x_1\}\in T_X\) is a direct successor of \(v_0\) in the representing tree \(T_X\), we have
\begin{equation}\label{eq:star6}
0= l_X(\{x_1\})<l_X(v_0)
\end{equation}
by definition of $T_X$.

Inequality \eqref{eq:star1} follows from \eqref{a2}, \eqref{eq:star5} and \eqref{eq:star6}.




Thus condition $(i)$ of Theorem~\ref{edftt} also is valid for the labeled rooted tree \(T^*(r^*,l^*)\). Consequently, by condition $(ii)$ of this theorem, we can find an ultrametric space \((Y,\rho)\) such that \(T^*(r^*,l^*)\) is isomorphic to the representing tree $T_Y=T_Y(r_Y,l_Y)$ of the space $(Y,\rho)$.


Let us prove equality \eqref{ned} for this $(Y,\rho)$. 


It follows directly from the 
 definition of the representing tree $T_X$ that $\{\{x\}: x \in X\}$ is the set of leaves of this tree. 
Similarly the set $\{\{y\}: y \in Y\}$ is the set of leaves of the representing tree $T_Y$ of the space $(Y,\rho)$. 
Moreover, equalities \eqref{avt1}, \eqref{avt2} imply that the set
\[
\{v^*\} \cup \{\{x\}: x \in X\}
\]
  is the set of leaves of the tree $T^{*}$.
It was shown above that $T^{*}(r^{*},l^{*})$ and $T_Y$ are isomorphic. 
Consequently the sets of leaves of the trees $T_Y$ and $T^{*}(r^{*},l^{*})$ have one and the same cardinality.
Now using the equality
\[
\{v^*\} \cap \{\{x\}: x \in X\}=\varnothing,
\]
we obtain
\[
|Y|
=
\left|\{\{y\}: y\in Y\}\right|
=
\left|\{v^*\}\cup \{\{x\}: x\in X\}\right|
=
\left|\{v^*\}\right|
+
\left|\{\{x\}: x\in X\}\right|
=
1+|X|.
\]
Equality \eqref{ned} follows.

Let us prove equality \eqref{gaf1}.





By Proposition \ref{ugrdd} the equalities 
\begin{equation}\label{gaf2}
C(X) = \bigcap_{p \in X} D_p(X) 
\end{equation}
and
\begin{equation}\label{gaf3}
C(Y) = \bigcap_{p \in Y} D_p(Y) 
\end{equation}
hold when
\begin{equation}\label{gaf4}
 D_p(X) =\{d(x,p)\colon x\in X\}
\end{equation}
for $p\in X$, and
\begin{equation}\label{gaf5}
 D_p(Y) =\{\rho(y,p)\colon y\in Y\}
\end{equation}
for $p\in Y$.


Using Corollary \ref{sefugi} we can rewrite equalities \eqref{gaf4} and \eqref{gaf5} as
\begin{equation*}
D_p(X) = \{ l_X(u) : u \in V(P_{\{p\},X}) \} 
\end{equation*}
and
\begin{equation*}
D_p(Y) = \{ l_Y(u) : u \in V(P_{\{p\},Y}) \},
\end{equation*}
where \(P_{\{p\},X}\) is the path in \(T_X\) joining the leaf \(\{p\}\) with the root \(X\) and \(P_{\{p\},Y}\) is the path in \(T_Y\) joining \(\{p\}\) with \(Y\).

Thus \eqref{gaf2}, \eqref{gaf3} imply the equalities
\begin{equation}\label{j1}
C(X)=\bigcap_{x\in X}\{l_X(u):u\in V(P_{\{x\},X})\}
\end{equation}
and
\begin{equation}\label{j2}
C(Y)=\bigcap_{y\in Y}\{l_Y(u):u\in V(P_{\{y\},y})\}.
\end{equation}

The set $\{\{y\}, y\in Y\}$ is the set of leaves of the representing tree $T_Y$ of the space $(Y,\rho)$. 


Since the trees $T^*(r^*,l^*)$ and $T_Y(r_Y,l_Y)$ are isomorphic as labeled rooted trees,
equality \eqref{j2} gives us the equality
\begin{equation}\label{j3}
C(Y)=\bigcap_{v\in L(T^*)}\{l^{*}(u):u\in V(P_{u,r^*})\},
\end{equation}
where $L(T^*)$ is the set of all leaves of $T^*$ and
$P_{u,r^*}$ is the path in $T^*$ joining the leaf $v\in L(T^*)$
with the root $r^*$ of $T^*$.
It follows from \eqref{avt1}, \eqref{avt2} that  the equality
\begin{equation}\label{j4}
L(T^*)=\{v^*\}\cup\{\{x\}:x\in X\}
\end{equation}
holds because $\{\{x\}:x\in X\}$ is the set of all leaves of $T_X$.

Now \eqref{j1}, \eqref{j3} and \eqref{j4} imply
\begin{equation}\label{j5}
C(Y)=C(X)\cap\{l^{*}(u):u\in V(P_{v^*,r^*})\}.
\end{equation}










Since $T^*(r^*,l^*)$ and $T_Y(r_Y,l_Y)$ are isomorphic and the leaves $\{x_1\}$, $v^* \in L(T^*)$ are direct successors of $v_0 \in V(T^*)$,
the equality
\begin{equation}
\{\, l^*(u) : u \in V(P_{v^*,r^*}) \,\}
=
\{\, l^*(u) : u \in V(P_{\{x_1\}},r^*) \,\}
\label{eq:star1-1}
\end{equation}
 holds by Lemma \ref{sabra}.
The definition of $T^*(r^*,l^*)$ implies the equality
\begin{equation}
\{\, l^*(u) : u \in V(P_{\{x_1\},r^*}) \,\}
=
\{\, l_X(u) : u \in V(P_{\{x_1\},X}) \,\},
\label{eq:star2-1}
\end{equation}
where $P_{\{x_1\},X}$ is the path in $T_X$ joining the leaf $\{x_1\}$ with the root $X$.
Equalities \eqref{j1}, \eqref{eq:star1-1} and \eqref{eq:star2-1} give us the inclusion
\begin{equation}
\{\, l^*(u) : u \in V(P_{v^*,r^*}) \,\} \supseteq C(X).
\label{eq:inclusion}
\end{equation}
The last inclusion together with equality \eqref{j5} gives us equality \eqref{gaf1} as required.

The proof is completed.
\end{proof}

\end{document}
